\section{bruitage et cryptage de life}

Il est assez facile de reconnaitre un point fixe, un cycle limite, le
chaos\ldots On fait des mesure, on cherche des divergences\ldots
Maintenant, je peu produire~(je ne sais rien de leur densit\'e
d'existence) une r\`egle universelle, \arule{Life} par exemple, qui aura
l'apparence d'une r\`egle cahotique\footnote{%
  Je ne peut pas lui donner l'apparence d'une r\`egle point fixe
  evidemment, mais je peu lui donner l'apparence d'une r\`egle cahotique
  si je veux} %
face \`a un equipement qui cherche avec des mesures statistiques.
%%
C'est simple, je rajoute des plan bruit\'es,
et la moyenne des comportement sera un comportement bruit\'e.
%%
Il y a un plan qui en fait calcule life\footnote{%
  Possibilit\'e d'un crytpage suppl\'ementaire de premier niveau?
  Peut \`etre par un calcul reversible.}. %
La regle life est cacul\'e sur deux plan, et les cellules inutilis\'e sur
un plan seront brouill\'e. %
Le mouvement de bit sur un ensemble de plan semble al\'eatoire, mais si
je regarde tel bit dans tel plan et\ldots~(si je fais une
transformation) je retrouve life. %
A~impl\'ementer:~Le cas le plus simple: 3 plan de bits, 2 qui font
rien, maximalement cahotique, et un qui fait life. %
Si je consid\`ere ma r\`egle comme une r\`egle 3 bit \`a 8 \'etats je vais dire
cette r\`egle la est cahotique, je ne peu rien lui faire calculer, je
sais pas comment faire. %
Si on m'apprend quil faut regader un seul bit, qui n'est pas ramenable
\`a un \'etat, comme ca, on a life. %
Alors la question, quel est le pourcentage des r\`egles en apparence
chaotique dans lequelles on pourrais trouver quelque chose qui
appartient plus \`a la classe 3 mais \`a la classe 4.

\section{probl\`eme de plongement}

Exemple: je pars des automate \`a 2 dimensions,
je prends le voisinnage de von-neuman. %
J'ais $2^{32}$ r\`egle spossibles. %
Je passe au voisinnage de Moore,
j'ais $2^{512}$ r\`egles possibbles. %
Parmi les $2^{512}$ r\`egles de Moore je trouve les $2^{32}$ du voisinnage
de von-neuman. %
Ca semble etre une petite partie. %
C'est un peu plus compliqu\'e que ca. %
Parceque je dois probablement les prendre sur la diagonales. %
C'est \`a dire faire une rotation \`a 45 degr\'e sur la table
et dire je les ai au mois deux fois. %
Je touve aussi tous les automates 1D sur les lignes les colonnes et les diagonale. %
Y sont pas nombreux non plu, mais en m\^eme temps on les recombine dans l'espace. %
Evidemment si je passe en trois dimension je vais retrouver toutes les
familles 2D et 1D par faisceaux dans tous les sens. %
Alors \`a chaque fois, a prioris le rapport entre le nombre de r\`egles
possible quand on augment la dimentionalit\'e ou le voisinnage, ou le
rayon, \emph{et} le nombre de r\`egle qui se trouve enchass\'e plong\'e dans
l'espace de dimmension sup\'erieure\ldots On peut peut-\^etre consid\'eere
qu'il est n\'egligeable, mais en tout cas il permet d'expliquer quelque
chose, c'est que si j'ais une r\`egle universelle en voisinnage de
von~Neumann, elle est aussi en voisinage de Moore elle est \`a des
endrois dans l'espace des r\`egles qui font quon n'aura pas un lieu bien
determin\'e de la pr\'esence des r\`egles presentant certaines
caract\'eristique puisque on a toute celle des espaces d'avant,
celle qui ne tiennent pas compte de certain voisinnage. %
(rapport \bialt{dimension}{nombre d'etat} pour compter nombre d'enchass\'e).
Essayer de formaliser.

\section{Tactique sociale}

Je suis poli avec mes voisins mais si l'un d'entre eux manifeste \`a mon
\'egard autre chose que cette bienveillante indiff\'erence\footnote{%
  S'il fait mine de se m\'eler de mes affaires, par exemple.}, %
je passe en mode agressif~(verbal), le plus rapidement, clairement et
brutalement possible.

Les tactiques pouvant \^etre consid\'er\'e comme des algorithmes\footnote{%
  Et les strat\'egies comme des heuristiques.}, %
Cette histoire peut inspirer une fonction de transition.

Si l'on ajoute des \`especes ont peut se demander quelles sont les
conditions d'apparition d'alliances.

\section{short notes}

\begin{itemize}

\item Incertitude, Heisenberg, unit\'e multitude. %
  Deterministe localement, non deterministe globalement. %
  Inverse du monde des particules~(monde quantique)

\item%
  Comparaison de $\lambda$ sur un petit espace de r\`egles~(inf\'erieure \`a
  $2^{16}$) avec lambda sur l'ensemble des r\`egles additives, le sous
  espace des r\`egle additives? %
  La distribution de lambda non lin\'eaire explique le param\`etre critique,
  la majorit\'e des r\`egles est dans la r\'egion critique?
  Pattern/s\'equences, espace/temps. %
  L'espace se voi, le temps s'entends, ou s'\'ecoute. %
  Synth\`ese picturale par AC. %
  Utilisation de la synthese sonore pour la repr\'esentation des s\'equences~(ask vi). %
  comparer distribution de lambda sur un \'echatillon al\'eatoire et sur un
  subset determin\'e~(comme les r\`egles additives)
  par rapport \`a la distribution sur un ensemble complet. %
  Ac de surfaces, interfaces entre deux milieux. %
  Tridimensionel mince. %
  Couches de 4 \`a 32 cellules d'epaisseur. %
  Interface eau/air, liquide/gas, vague, propagation d'ondes de surface, sable. %
  Les tas de sables, burst de chaos dans plage d'equilibre,
  mod\`ele cellulaire des tas de sable. %
  Les pyramides \`a degr\'e comme des tas de sable. %
  
\item \`a propos de l'enchassement des espaces de r\`egles. %
  L'hypoth\`ese ``computation at the edge of chaos'' postule une
  transition de phase abrupte.

\item Cr\'eation d'un catalogue de de r\`egles pouvant constituer un
  niveau d'instruction \'el\'ementaire pour une machine cellulaire.

\item%
  $$O(\lambda) = 2^N$$
  $$Space(\lambda) = \log_2(N)$$
  $$O(E) = N 2^N$$
  $$Space(E) = \log_2(\frac{N}{2}) N$$
  
  interpolation de $\lambda$ croissnt de 0 \`a N.
  interpolation de E de 0 \`a 0.
  par $\frac{N}{2}$ pour $x = \frac{2^N}{2}$?

\item%
  3 solutions d'impl\'ementation:

  \begin{itemize}
  \item filtre, langage, batch
  \item shell
  \item click and point
  \end{itemize}

  langage = lex, yacc matlab - txl; shell = sh, tcl; %
  tcl = langage; tcl + tk = click and point; %
  methodologie objet, tcl + itcl = objet; athena widget = trop fig\'e,
  trop lourd. %
  Mod\`ele click an point = khoros. %
  Alternative: objet c ou c++, tcl et widget tk; %
  shell = devellopement; %
  langage = maturation; visualisation = something.
  
\item%
  R\`egles additives, 10 entr\'ees dans la table secondaire,
  4 \'etats par entr\'ee, z\'ero, un unchanged et flip. %
  Somme sur les 8 voisins = 9 entr\'ee; %
  somme sur les 8 voisins + le  centre = 10 entr\'ees.
  
  Pour 10 entr\'ees: $4^{10}$ ou $2^{20}$ ou $\approx$ 1 millions de
  r\`egles possible.
  
\item%
  $\lambda$ + E pour la cellule centrale + E pour les voisinnes
  
  E petit pour $\lambda$ petit
  E petit pour $\lambda$ grand
  
  au moins sur un exemple E = $\lambda_{critique}$, pour $\lambda$ = 0.5.
  
  question: variance de $\lambda$ = $\lambda_{critique}$?
  
  Calculer pour un espace ce r\`egle, les densit\'ees d'apparition des
  diff\'erentes valeurs de $\lambda$, et le nombre de valeur diff\'erentes
  de $\lambda$ et le nombre de r\`egles \`a la temp\'eratute critique par
  rapport au nombre total de r\`egles.

\item Complexe = en croissance structurelle?

\item Combien y a til de r\`egles universelle dans nu espaces de r\`egles. %
  Le probl\`eme de l'universalit\'e d'une r\`egle est-il d\'ecidable?

\item Life est une r\`egle unique parceque de nombreuse personnes on
  travaill\'ee ensemble \`a la construction de la preuve de son
  universalit\'e.

\item%
  Question du cout de l'architecture de d\'evellopement. %
  Architecture mat\'erielle et logicielle syst\`eme. %
  Rapport entre la puissance des algorithmes cellulaires et le cout
  de l'architecture mat\'erielle. %
  La moins ch\`ere des machine haut de gamme. %
  30 mips, 30 m\'ega-bytes 1 giga disk, 1 million pixel, 256 couleur. %
  Relation entre les unit\'e cellulaires et les unit\'e d'espace
  et de vitesse des machines. %
  Univers de r\`egles, univers cellulaires accessible pour une puissance
  machines donn\'ee. %
  Capacit\'e machines demain: 300 mips, 256 m\'ega. %
  Notre architecture de d\'evellopement est l'architecture de r\'ef\'erence. %

\item les AC et la complexit\'e par le bruit? %

\item Pour faire un univer: l'espace, le temps, la loi~(de transition
  local) et une rupture de sym\'etrie sur l'espace pour permettre \`a la
  loi de s'exprimer. %

\item Tas de sable, equilibre ponctu\'e, applicable \`a quoi d'autre?

\item Comportement cellulaire global cyclique, comment?  combien de
  plan de bits sont necessaire~(a par solution triviale de compteur local). %

\item = plan de coupe = trance, slice, ligne de coupe = carotte,
  morceau = slab. %

\item \'Enumeration des r\`egles par types pour un espace complet. %

\item totry:~The history of the density of the run time is the sum of the
  density of 1 thru 8 connected cell. %
  It is expected to find a correlation between AC class and distribution
  of various connected sum. %

\item L'histogrammes des densit\'e locales permet de calculer la taille des
  fronti\`ere dans anneal par exemple. %

\item A reference to a stochatic CA being assigned that a
  deterministic CA if the random number generator is part of the rule
  can be found in latice gas method for partial differential equation
  a CA for burger equations p 486. %
  They use wolfram ca random number generator, or random bit generator
  from unbiased distribution and using logical expression logical
  operation to get biased distribution. %
  In that paper on can find a expression for the covariance of a CA.

\item Wolfram origin of randomness in physical system: ``if in fact, no
  polynomial time procedure ca detect patterns in the sequence then the
  sequence can be considerer effectively random for practical purposes.''

\item Pour chaque r\`egle exp\'eriment\'e, rechercher la divergence pour un bit de
  variation de la condition initiale.

\item Majority with flip, anneal with flip.

\item Try to propagate constraint, that number of change from 0 to 1
  or 1 to 0 always stay equal.

\item \todo{Definir ergodicity}. %
  Trouv\'e une note dans \cite{Queau89}:~Un syst\`eme est dit ``ergodique''
  si la moyenne temporelle de ses sorties \'equivaut \`a la moyenne des
  sorties d'un grand nombres de syst\`emes fonctionnant simultan\'ement.

\item In prospect for a lattice gas computer, in latice gas method for
  partial differential equation:\\
  %%
  ``the question about 24 bit model for 3 dimensional latice gas is: how
  can this large number of rule can be reduce by factoring or grouping
  them. %
  Reduce the size of the rule representation in the lookup table. %
  What is the fundamental dimension
  of the rule set?''
  %%
  Essayer dimplementer la table du mod\`ele 24 bit calculer la valeur
  de~$E$, essayer de d\'ecouper en plusieurs tables par r\'eduction via
  un~DFA.

\item Quelle est l'intersection entre un espace de r\`egles et la
  ph\'enom\'enologie naturelle?

\item Same book page 214:\\
  %%
  g\'en\'eraliser la question 7b au automates cellulaires.
  %%
  page 215:\\
  %%
  ``A different type of test for latice gas algorithm was due to be
  important as well. %
  One should perform the standard num\'erical test of increasing the grid
  resolution while maintaining fixed all the parameters describing the
  problem the goal would be to check that the higher resolution results
  is identical to those obtain with the lower
  numerical resolution.''
  %%
  Generaliser \`a tous les AC, et essayer appliquer sur anneal par
  exemple.

\item Utiliser anneal pour des experiences du m\^eme type que celle des
  latice~gas. %
  C'est \`a dire avec des source ou des flux.
  
\item Tester le changement d'echelle pr\'ealablement d\'ecrit sur anneal dans
  ces conditions. %
  Cad.\ pe.\ une source fixe de 1 de taille proprortionelle identique
  pour des nombre de cellulles totale different et les diff\'erence
  produite.

\item Introduire les probl\`eme de la visualisation \`a partir de l'article en
  question page 215.

\end{itemize}

%%
%% Local Variables:
%% mode: auto-fill
%% iso-accents-minor-mode: nil
%% TeX-master: "top"
%% TeX-command-default: "mytex"
%% Local IspellParsing: "latex-mode ~latin1"
%% Local IspellPersDict: "~/.ispell_lisp"
%% LocalWords: "lookup"
%% End:
